% ============================================================================
% NEW MATERIAL, October 6, 2026 (agent draft; not in JMP_DS_draft).
% Results section skeleton. Every number is from the base point
% output/model/production/2007 (Oct 6 refit, loss 15.07, provisional) unless a
% comment says otherwise:
%   untargeted moments: output/model/production/2007/plotted_long.csv
%   price and credit arms: output/model/jmp_draft_deck_20261005/
%     refit_best_15p07_credit/results.csv (fixed price fp_*, stationary GE ge_*)
% Results that exist only at chain 11, chain 13 or the Sept 14 point are
% NOT typeset; they are listed in comments as the rerun queue.
% ============================================================================
\providecommand{\mocktodo}[1]{\textcolor{red}{[TODO: #1]}}

\section{Results}
\label{sec:results}

This section uses the calibrated economy to ask how housing costs and access
to credit shape fertility. I first evaluate the model on moments it was not
asked to match. I then raise the price of housing and, separately, relax the
down-payment requirement, and trace each change to births, ownership and the
size of homes. Finally, I describe the historical exercise, in which the
economy moves from 2007 to 2023.

\subsection{Moments not targeted}

The model reproduces the distribution of first births by age, although only
the mean age was targeted: in the model 34.6 percent of first births occur
before age 22 and 11.9 percent at ages 30 to 33, against 34.2 and 14.3 percent
in the data. It misses in three places, each informative about the mechanism.
% [MODEL] plotted_long.csv first_birth_age, cells 18,22,...; data NCHS.

First, young households own too much. At age 24, 42.9 percent of households
in the model own their home, against 23.5 percent in the data; by age 40 the
two agree, at 67.2 and 66.5 percent. The model matches ownership at ages 30
to 55 by having households buy earlier than they do in the data.
\mocktodo{name the data source of the ownership-by-age profile in the
assessment packet (model\_data\_assessment.pdf).}

Second, the first child raises space too much. Households with one child
occupy 7.3 rooms in the model and 5.9 in the data, against 5.0 and 5.1
without children. This is the other side of the miss on the first-birth room
response in Table~\ref{tab:quant-fit-oct06}: the model adds too few rooms
around the birth itself but ends up with too many once the child is there.
\mocktodo{these two statements are hard to reconcile without the event-time
paths; check the rooms-by-children observer (cross-section, all ages) against
the event-study observer before keeping the sentence.}

Third, and most important, the model gets the income gradient of fertility
wrong. Among households aged 24 to 26, those in the bottom third of income
have 0.15 children in the model and 0.61 in the data, and those in the top
third 0.94 and 0.23. By ages 40 to 44 the model's gradient is still positive,
1.50 against 2.31, while the data's is slightly negative. In the model, the
cost of the parent floor is a large share of income for the poorest
households, so they postpone or forgo children; in the data they have them
earliest. Section~\ref{sec:results-price} shows why this matters for the
response of births to housing costs.
% [MODEL] plotted_long.csv income_24_26 (Model .149/.554/.943; Data
% .608/.453/.226), income_40_44 (Model 1.498/1.969/2.311; Data
% 2.013/1.891/1.851). The data columns differ from Table 1 (0.596/.../1.784):
% different grouping (assessment packet vs measurement/tables.json).
\mocktodo{use one definition of the CPS income thirds in Table~\ref{tab:emp-fertility-income}
and here; the assessment packet and the Section 2 table differ by up to 0.28
children at 40--44.}

\subsection{Housing costs and births}
\label{sec:results-price}

The cost of a child in the model has two parts. The equivalence scale raises
the spending needed to keep consumption per person constant, and the parent
floor requires $h_P$ rooms before housing yields any services. The second part
is large. At the calibrated price, a room rents for 0.141 per period, so the
floor costs $0.141\times2.585=0.363$ per period, or 9.1 percent of mean annual
earnings each year. For a household entering at 18 with average productivity
it is 12.6 percent of its gross earnings.
% [MODEL] r = user_cost_rate 0.18028 x P 0.77952 = 0.14053 per room per
% period; h_P 2.58534; annual 0.0908; entrant annual gross 0.72055.

I raise the house price and the rent by 10 percent and hold all other prices
and policies fixed. Births fall by 4.9 percent. Ownership at ages 18 to 29
falls by 2.8 percentage points, and young households occupy 6.6 percent fewer
rooms. Table~\ref{tab:results-experiments} reports the full set of
responses.
% [MODEL] fp_price110: births -4.8828%, own 18-29 -2.7936 pp, rooms 18-29
% -6.5557%, mean rooms -5.9773%.
\mocktodo{split the price effect into an income effect and the cost of the
floor at the base point (rerun the Shapley split; it exists only at chain 11,
where the floor accounts for 60 percent and income for 45 percent).}
\mocktodo{add the stationary general-equilibrium version of the price
experiment; only the fixed-price arm exists at the base point.}

\subsection{Credit and births}
\label{sec:results-credit}

Because rent equals the user cost of an owner and households and landlords
face the same interest rate, a room costs the same per period whether it is
rented or owned. The down-payment requirement therefore governs when a
household can buy, not what its space costs over the life cycle. I raise the
financed share $\phi$ from 0.80 to 0.95, holding prices fixed. Ownership at
ages 18 to 29 rises by 25 percentage points, from 43 to 68 percent, and
births fall by 1.4 percent. With $\phi=1$, ownership rises by 29 points and
births fall by 1.3 percent.
% [MODEL] fp_phi095: births -1.4400%, own 18-29 +25.10 pp (0.4335 -> 0.6845);
% fp_phi100: births -1.3480%, own 18-29 +28.82 pp.

The fall in births comes from the homes that newly eligible buyers choose.
Among childless households aged 18 to 37, the share that owns a two-room home
rises from 6.7 to 42.2 percent. Two rooms are below the parent floor, so a
household in such a home must move before it can have a child, and moving
costs 6 percent of the home's value out of little equity.
% [MODEL] refit_best_15p07_credit/README.md starter-home check: renter /
% 2-room owner / larger owner 74.1/6.7/19.3 -> 28.4/42.2/29.4 at phi .95.
\mocktodo{decompose the credit effect at the base point into the change in
decisions at a given state and the change in the distribution of states
(exists only at chains 11 and 13; at chain 11 the distribution term dominates).}

In a stationary equilibrium, completed fertility must equal replacement, so a
change in credit cannot change births per household in the long run; it
changes the price at which the population replaces itself, and the size of the
population. With $\phi=0.95$, the stationary price is 2.9 percent lower and
the number of households 3.0 percent lower.
% [MODEL] ge_phi095: price -2.936%, population_scale 0.97053, births
% per household -0.08% (closure); ge_phi100: price -1.978%, scale 0.98365.
\mocktodo{check the sign and wording of the population result with Tommaso
before keeping it; it is a stationary comparison, not a transition.}

\begin{table}[t]
\centering
\caption{Responses to housing costs and credit, 2007 steady state}
\label{tab:results-experiments}
\small
\begin{tabular}{lrrrr}
\toprule
 & Births (\%) & Ownership 18--29 (pp) & Rooms 18--29 (\%) & Price (\%) \\
\midrule
\emph{Fixed prices} & & & & \\
House price and rent $+10\%$ & $-4.88$ & $-2.79$ & $-6.56$ & $+10.0$ \\
Financed share $0.80\to0.95$ & $-1.44$ & $+25.10$ & $-0.84$ & 0 \\
Financed share $0.80\to1.00$ & $-1.35$ & $+28.82$ & $-0.71$ & 0 \\
\addlinespace
\emph{Stationary equilibrium} & & & & \\
Financed share $0.80\to0.95$ & $-0.08$ & $+27.03$ & $+1.32$ & $-2.94$ \\
Financed share $0.80\to1.00$ & $-0.10$ & $+28.92$ & $+0.97$ & $-1.98$ \\
\bottomrule
\end{tabular}
\par\smallskip
\begin{minipage}{0.92\linewidth}
\footnotesize\textit{Notes:} Changes relative to the 2007 steady state.
Births are births per household per period. In the stationary equilibrium the
price solves the replacement condition, so births per household change only
through numerical tolerance. The property-tax rebate is rebalanced in every
case.
% [MODEL] results.csv rows fp_price110, fp_phi095, fp_phi100, ge_phi095,
% ge_phi100.
\end{minipage}
\end{table}

\subsection{The fertility decline, 2007--2023}
\label{sec:results-transition}

\mocktodo{Not yet run at the base point. The planned exercise starts from the
2007 steady state and lowers the child benefit $\psi$ by two unanticipated
permanent changes, in 2007 and 2015, fitting completed fertility of 1.861 in
2012--2015 and 1.646 in 2020--2023 (CALIBRATION\_STATUS.md, outstanding item
1). Figures to report, following the September 14 presentation: the fit of
fertility to the data, and the four-panel path of fertility, households,
housing and the house price. The existing transition figures come from an
earlier, unconverged calibration point and are not used here.}

\subsection{A housing policy}
\label{sec:results-policy}

\mocktodo{The property-tax experiment of the September 14 presentation
(doubling the tax with an equal rebate, from the 2023 economy) exists only at
the earlier point: births $+1.0$ percent after forty years. Rerun from the
base point's 2023 economy once the transition is accepted.}
